% GDZ PPN236018647,native573lower–574;HKUST574–575.
\NumberedItem{102}
\indent Medium Arithmetico-Geometricum ipsum est quantitas integralis. Dem[onstratum].
\par\noindent\hfill[1799] Dec. 23.

\medskip\begingroup\small
\indent Man könnte diese Notiz mit der Gleichung [20], oben S. 187, in Verbindung bringen, die ja vermöge der Integraldarstellung [16] von \(Q\) auch eine Integraldarstellung für den reziproken Wert des agM. liefert; für eine solche Verbindung spräche auch das hier wie dort später hinzugeschriebene demonstr[atum]. Dagegen wäre es nicht recht verständlich, weshalb \textsc{Gauss} neun Tage (14.—23. Dezember) gebraucht haben sollte, um aus der Darstellung [16] die Folgerung [20] zu ziehen; auch machen die Aufzeichnungen auf S. 11 der Scheda Ac dem Ansehen der Schriftzüge nach den Eindruck, unmittelbar hintereinander geschrieben zu
\par\endgroup

\clearpage\SourceLeaf{574}\ArithmeticHead{545}{R}{1799 SEPT.—1800 APR. 27.}
\begingroup\small
\noindent sein. Und in der Tat scheint die fast wörtliche Übereinstimmung mit der Bemerkung im art. 8. der Abhandlung von 1800, Werke III, S. 370: »hoc itaque modo media nostra arithmetico-geometrica ad quantitates integrales revocata sunt« eher darauf hinzudeuten, daß es sich hier darum handelt, daß der reziproke Wert des agM. als das Integral einer einfachen Differentialgleichung definiert werden kann, so daß der Abschnitt [III.], oben S. 181, hier als Beleg heranzuziehen wäre.
\par\noindent\hfill\textsc{Klein.\quad Schlesinger.}
\par\endgroup

\NumberedItem{103}
\indent In theoria formarum trinariarum formas reductas assignare contigit.
\par\noindent\hfill1800 Febr. 13.

\medskip\begingroup\small
\indent Dasselbe Datum Werke I, S. 476, Bemerkung zum art. 272. der \textit{Disquisitiones arithmeticae}, Werke I, S. 307. Während die Notizen Nr. 100 bis Nr. 102 mit den Aufzeichnungen auf S. 7—15 der Scheda Ac (siehe oben S. 184—193) in Verbindung stehen, gehört zu dieser Notiz eine Eintragung auf S. 22 der Scheda, die Werke II, S. 311 abgedruckt ist.
\par\noindent\hfill\textsc{Klein.\quad Schlesinger.}
\par\endgroup

\NumberedItem{104}
\indent Series\DiaryBracketNote{*}{Die Handschrift hat \textls[120]{Seriem}.}
\[
a\cos A+a'\cos(A+\varphi)+a''\cos(A+2\varphi)+\text{etc.}
\]
\noindent ad limitem convergit, si \(a,a',a''\) etc. constituunt progressionem sine mutatione signi ad \(0\) continuo convergentem. Demonstratum.
\par\noindent\hfill Brunov[ici, 1800] Apr. 27.

\medskip\begingroup\small
\renewcommand{\thefootnote}{**)}
\indent Es ist hier der Satz ausgesprochen, daß die Reihen
\[
\begin{aligned}
a+a'\cos\varphi+a''\cos2\varphi+\cdots,\\
a'\sin\varphi+a''\sin2\varphi+\cdots
\end{aligned}
\]
\noindent konvergieren, wenn die \(a,a',a'',\ldots\) \textls[120]{monoton}\footnote{Das liegt in dem \textls[120]{continuo} des Tagebuchtextes. Ohne diese Einschränkung ist der Satz bekanntlich nicht richtig.} der Null zustreben. Dieser Satz ist zuerst veröffentlicht und bewiesen von \textsc{Hjalmar Holmgren}, Journal de Mathématiques \textbf{16}, 1851, S. 186, mit einer überflüssigen Einschränkung schon von \textsc{C. J. Malmsten}, Nova acta Upsal. \textbf{12}, 1844, S. 255.
\par\noindent\hfill\textsc{Schlesinger.}
\par\endgroup
